\section{Schwarz maps on direct sums of matrix algebras}


% This is the finite-sum complete-positivity implication used for the
% transported composites in CPSV16, Appendix C.4, lines 1974--1995.  It does
% not establish Kraus representations for those transported maps; that is the
% tensor-specific preparation and partial-trace step tracked in issue #4438.


\begin{theorem}[Paired component maps for products of matrix algebras]
    \label{thm:mps_component_map_bijective}
    \lean{MPSTensor.ringEquiv_maps_single_support}
    \lean{MPSTensor.componentMap_surjective}
    \lean{MPSTensor.componentMap_bijective}
    \leanok

    For a ring equivalence between products of full matrix algebras, the map
    induced between each pair of matched components is bijective.
\end{theorem}

\begin{proof}\leanok

    Apply Theorem~\cite[``Paired component maps are bijections'']{QICLean2026Blueprint} to the matched block
    ideals.
\end{proof}

\begin{theorem}[Dimension preservation for matched matrix summands]
    \label{thm:mps_dim_preserved}
    \lean{MPSTensor.dim_preserved}
    \leanok
    An algebra automorphism of a product of full matrix algebras preserves
    the dimension of every pair of summands matched by its induced
    permutation.
\end{theorem}

\begin{proof}\leanok
    \uses{thm:mps_component_map_bijective}
    Restrict the automorphism to the two matched summands. The resulting
    component map is a bijective complex-linear map by
    Theorem~\ref{thm:mps_component_map_bijective}. Hence the two matrix
    spaces have equal complex dimensions, so their sizes have equal squares
    and therefore are equal.
\end{proof}


\begin{theorem}[Transport between spatial compressions]%
    \label{thm:corner_compression_partial_isometry_transport}
    \lean{Matrix.exists_partial_isometry_transporting_compressions}
    \leanok
    Let $P,Q\in\MN D$ be matrices, and suppose that
    $V_P,V_Q:\mathbb C^n\to\mathbb C^D$ satisfy
    \[
        V_P^*V_P=I,\qquad V_PV_P^*=P,\qquad
        V_Q^*V_Q=I,\qquad V_QV_Q^*=Q.
    \]
    Write $\Phi_P(X)=V_PXV_P^*$ and $\Phi_Q(X)=V_QXV_Q^*$ for the
    corresponding linear isomorphisms onto the two corners. For every unitary
    $W\in\MN n$, there is a matrix $U\in\MN D$ such that
    \[
        U=PUQ,\qquad U^*U=Q,\qquad UU^*=P,
    \]
    and, for every $X\in\MN n$,
    \[
        \Phi_P(WXW^*)=U\Phi_Q(X)U^*.
    \]
\end{theorem}

\begin{proof}\leanok
    Take $U=V_PWV_Q^*$. The four isometry identities give the three support
    identities by direct multiplication. Substitution of the two formulas
    for $\Phi_P$ and $\Phi_Q$ gives
    \[
        U\Phi_Q(X)U^*
        =V_PW(V_Q^*V_Q)X(V_Q^*V_Q)W^*V_P^*
        =\Phi_P(WXW^*).
    \]
\end{proof}

\begin{theorem}[Spatial implementation of an isomorphism between corners]%
    \label{thm:corner_star_isomorphism_partial_isometry}
    \lean{Matrix.exists_partial_isometry_implementing_corner_linearEquiv}
    \leanok
    Let $P,Q\in\MN D$ be orthogonal projections. If
    \[
        \Psi:Q\MN DQ\longrightarrow P\MN DP
    \]
    is a bijective complex-linear map preserving multiplication and the
    adjoint, then there is a matrix $U\in\MN D$ such that
    \[
        U=PUQ,\qquad U^*U=Q,\qquad UU^*=P,
    \]
    and
    \[
        \Psi(X)=UXU^*
    \]
    for every $X\in Q\MN DQ$.
\end{theorem}

\begin{proof}\leanok
    \uses{thm:corner_compression_partial_isometry_transport}
    Choose spatial compressions
    \[
        \Phi_Q:\MN{n_Q}\simeq Q\MN DQ,\qquad
        \Phi_P:\MN{n_P}\simeq P\MN DP.
    \]
    The composite $\Phi_P^{-1}\Psi\Phi_Q$ is a linear isomorphism, so equality
    of dimensions gives $n_Q^2=n_P^2$, and hence $n_Q=n_P$. If this common
    rank is zero, both projections vanish and the conclusion holds with
    $U=0$. Otherwise the composite is a star-algebra automorphism of
    $\MN{n_Q}$. The unitary Skolem--Noether theorem writes it as
    $X\mapsto WXW^*$ for a unitary $W$. Applying
    Theorem~\ref{thm:corner_compression_partial_isometry_transport} to $W$
    gives the required matrix $U$.
\end{proof}
